\section{Conclusão}

Este trabalho avaliou a ConvNeXt-Tiny na triagem binária de células cervicais
da CRIC Cervix usando validação cruzada em cinco folds agrupada por imagem. O
modelo apresentou AUC ROC de $0{,}9709\pm0{,}0059$ e AUC PR de
$0{,}9607\pm0{,}0121$. A comparação complementar com ResNet-50 nos mesmos folds
mostrou desempenho operacional muito próximo, mas menor AUC ROC e AUC PR
médias para a ResNet-50. As diferenças em sensibilidade, especificidade e
acurácia balanceada foram pequenas e variaram entre folds, indicando que a
composição das imagens e o limiar operacional também influenciam o desempenho
observado.

A análise reforça a importância de reportar média, desvio-padrão e variação
fold a fold em estudos com recortes celulares. O particionamento por imagem
reduz vazamento entre treino e avaliação e torna a comparação entre ConvNeXt-Tiny
e ResNet-50 mais controlada, pois ambas as arquiteturas enfrentam os mesmos
subconjuntos de imagens. Ainda assim, como a CRIC não fornece identificadores de
paciente, o protocolo não garante independência no nível do paciente ou da
lâmina. Ademais, a ResNet-50 foi ajustada com orçamento de épocas menor, restrito
à janela em que a ConvNeXt-Tiny alcançou seus melhores checkpoints; embora isso
padronize o critério de seleção, uma comparação com orçamento idêntico permanece
como trabalho futuro.

Em relação à literatura, o resultado deve ser lido com cautela: acurácias
elevadas reportadas em Herlev, SIPaKMeD e outras bases não são diretamente
intercambiáveis com a CRIC, que contém citologia convencional com maior
variação visual, sobreposição celular e fundo heterogêneo. A contribuição deste
trabalho está menos em comparar números absolutos com bases distintas e mais em
estabelecer uma avaliação reprodutível, agrupada por imagem e pareada entre
arquiteturas.

Os resultados devem ser interpretados como evidência experimental para
priorização celular, não como validação clínica de um sistema diagnóstico.
Trabalhos futuros devem avaliar calibração, escolha de limiar, validação
externa, análise por categoria Bethesda, detecção em campos completos e
agregação no nível da lâmina. Código, configurações e materiais complementares:
\url{https://github.com/Kariellyy/Pesquisa-Karielly}.
